\section{Introduction}

An unmanned aerial vehicle (UAV) is an aircraft with no onboard pilot; in this thesis, UAV means a quadcopter, which
is underactuated, with four control inputs for six degrees of freedom \cite{kotaruDifferentialflatnessControlQuadrotors2018}.
Many UAV tasks, from delivery to construction, require carrying a payload \cite{kardaszDronesPossibilitiesTheir2016,
olleroPresentFutureAerial2022, thusooQUADROTORSPRESENTERA2021}, and one UAV's thrust limits what it can lift while a
suspended load swings and costs control effort to reject \cite{geronelDynamicResponsesDue2023}. Several UAVs carrying
one load through cables raise the liftable mass, constrain its swing and can control its full pose
\cite{tudelft, mohammadiDecentralizedMotionControl2018}, at the price of keeping every cable taut and sharing the load
so that no vehicle is overloaded \cite{kotaruDynamicsControlQuadrotor2017}. A large body of work addresses tracking,
tension allocation and disturbance rejection for fixed teams, mostly in simulation \cite{lv2026review,
estevezReviewAerialTransportation2024}.

Changes in team composition during a mission have received much less attention. In this thesis, reconfiguration means
a change of team membership, a UAV leaving or joining while the payload is carried, as distinct from a fixed team
changing its formation \cite{jackson2020scalable}. A UAV leaving a carried load has been flown
\cite{quan2025swarm, zeng2025marl}, but no hardware work read in full gives a number for the load's attitude through
the event, and no work has shown a UAV attaching to a payload that a team is already carrying (Section~\ref{sec:lit_team}).
A UAV may need to leave because of a low battery, a fault or a new task \cite{farrasFaultTolerantControlMultiQuadcopter2024},
and a replacement that joins in flight lets the mission continue without landing the team; a mission could also start
with the fewest UAVs and add more only when needed (Figure~\ref{fig:mission_profile}).

\tikzset{
  block/.style={
    rectangle, rounded corners=4pt,
    minimum width=5.3cm, minimum height=0.85cm,
    draw, thick, align=center,
    font=\small
  },
  transport/.style={block, fill=blue!6, draw=blue!40},
  event/.style={block, fill=orange!12, draw=orange!60},
  docking/.style={block, fill=purple!8, draw=purple!45},
  arrow/.style={-Stealth, thick},
  phase/.style={font=\footnotesize\color{gray}}
}

\begin{figure}[h]
\centering
\begin{tikzpicture}[node distance=0.6cm and 1.2cm]

\node (p1) [transport]
{\textbf{Initial transport}\\[-2pt]\scriptsize $N$ UAVs follow the planned trajectory};

\node (ev) [event, below=of p1]
{\textbf{UAV departure}\\[-2pt]\scriptsize failure / low battery};

\node (p2) [transport, below=of ev]
{\textbf{Reduced team}\\[-2pt]\scriptsize $N\!-\!1$ UAVs maintain stability};

\node[phase, left=0.4cm of p1] {Phase 1};
\node[phase, left=0.4cm of ev] {Phase 2};
\node[phase, left=0.4cm of p2] {Phase 3};

\draw[arrow] (p1) -- (ev);
\draw[arrow] (ev) -- (p2);

\node (p3) [docking, right=1.2cm of p1]
{\textbf{Replacement approach}\\[-2pt]\scriptsize localises docking point};

\node (p4) [docking, below=of p3]
{\textbf{Attachment}\\[-2pt]\scriptsize guided electromagnetic docking};

\node (p5) [transport, below=of p4]
{\textbf{Full team resumes}\\[-2pt]\scriptsize load redistributed to $N$ UAVs};

\node[phase, right=0.4cm of p3] {Phase 4};
\node[phase, right=0.4cm of p4] {Phase 5};
\node[phase, right=0.4cm of p5] {Phase 6};

\draw[arrow] (p3) -- (p4);
\draw[arrow] (p4) -- (p5);

\draw[arrow, thick] (p2.east) -- ++(0.8,0) |- (p3.west);

\end{tikzpicture}

\caption{Mission profile for dynamic UAV replacement during transport.}
\label{fig:mission_profile}
\end{figure}

This thesis addresses the control problem of keeping the payload stable through a departure and a mid-mission
addition, Phases 2, 4 and 6 of Figure~\ref{fig:mission_profile}; the docking mechanism of Phases 4 and 5 is the work
of a collaborating student. Validation is in simulation and on a hardware rig.
